\section{Conclusion}
\label{sec:conclusion}

Recursive AI-assisted systems are increasingly practical, but recursive execution alone is no longer the primary systems challenge. As recursion depth, artifact volume, and iteration speed increase, local progress can outpace global understanding. Under these conditions, recursive drift is not an edge case; it is a structural tendency. The central engineering question therefore shifts from how to execute faster to how to keep recursive workflows synchronized, observable, and governable as they evolve.

This paper has argued that \reflector{} should be understood as reflective synchronization infrastructure for that problem space. Its core contribution is not autonomous recursive control, and it is not presented as a replacement for human engineers. Instead, \reflector{} combines bounded recursive execution, reflective auditing, synchronization checkpoints, and governance layers into a coordination architecture that keeps continuation decisions tied to current evidence and explicit intent. In this framing, recursive orchestration is specification-driven and recursively observable: progress is treated as valid only when artifacts remain interpretable and reconcilable across loops.

The broader implication is human-centered. Mixed-initiative recursive systems work because humans remain embedded at the boundaries that determine trajectory: scoping delegation, interpreting ambiguity, authorizing continuation, and redirecting when drift emerges. AI systems can expand execution capacity, but synchronization practices preserve alignment between generated artifacts and human intent over time. Without that reflective coupling, recursive throughput can amplify incoherence as quickly as it amplifies output.

As recursive AI-assisted development matures, synchronization may become a foundational systems concern alongside correctness, performance, and reliability. Reflective synchronization architectures offer one grounded response: collaborative recursive infrastructure that prioritizes coherence, interpretability, and governance without over-claiming perfect control. The claim of this paper is therefore intentionally modest but consequential: the long-term viability of recursive systems will depend not only on what they can produce, but on whether they can remain synchronized with the humans who must understand, govern, and trust them.
